\section{实验矩阵、预测设计与验证方法}
\label{cr:design}
\subsection{固定矩阵与响应定义}
R50-main主设计包含四种栅长、两种粒子和两个方向，每个组合预定四个独立种子，
共64个独立粒子事件；每个事件分别施加于存储0和存储1，形成128个响应槽位。
配对设计有助于识别位值不对称，但同一沉积事件的两种初态不构成两个独立入射。
数值复跑、网格细化、参考诊断和氧化层配对同样不增加独立事件数。
所有种子与工况在读取正式响应前登记，失败不通过更换种子补成有利样本。

R50-main使用如下连续指标：
\begin{align}
\Delta Q&=C_{\mathrm{store}}\,[V_{\mathrm{SN,event}}(2.8\ns)
 -V_{\mathrm{SN,control}}(2.8\ns)],\\
D(t)&=V_{\mathrm{BL,signal}}(t)-V_{\mathrm{BL,neutral}}(t),\\
M&=(2b-1)D(4\ns),\qquad
\Delta M=M_{\mathrm{control}}-M_{\mathrm{event}}.
\label{cr:metrics}
\end{align}
其中control是同位值无粒子对照，neutral是中性参考，两者功能不同。
$\Delta Q$保留符号，负值表示事件SN低于对照；$\Delta M>0$表示裕量退化。
对存储0而言，节点进一步下降可能使判据裕量增加，此时$\Delta M$可以为负，
不应把所有变化都强制写成非负的“损伤”。

主判据取$\varepsilon=5\mV$：$M>\varepsilon$为正确，
$M<-\varepsilon$为判据误读，其余为临界。
同时保留$2.5/10\mV$阈值结果。接近边界而未完成规定数值复验、
或粗细离散分类不一致时，另列数值不确定；失败和缺失保持各自状态。
统计按冻结分母报告，不用可判子集替代预定样本。

\subsection{分组拟合与测试前冻结}
训练集取$L=0.18/0.35\um$，共32个事件、64个响应；
$0.25\um$为域内留出，$0.45\um$为尺寸外推，各16个事件、32个响应。
同一物理事件及其配对结果保持同组，避免相同源信息同时出现在训练和测试中。
预测使用真实Geant4沉积量，目标是在给定沉积与无粒子条件下近似TCAD响应。
它有望减少后续器件求解成本，但仍须获得粒子输运输入。

分别针对两种位值的$\Delta Q$与$\Delta M$建立四个岭回归模型，
特征为$Q_{\mathrm{dep}}$、$Q_{\mathrm{dep}}/L$和
$Q_{\mathrm{dep}}\sin\theta$。这些项以电荷注入规模、尺寸及方向为物理启发，
并不是已推导出的通用器件定律。与训练均值基准比较，可检查回归是否比
忽略输入变量更有信息；数据驱动单粒子响应建模的相关研究也为这种代理思路
提供背景\cite{xu2021}。

训练采用留一物理工况交叉验证（LOCO）。训练区有八个尺寸、粒子、角度组合，
每折将一个完整组合的四个重复留出；两位保持相同的分组边界。
标准化均值和方差只由本折训练数据计算。候选正则化参数为
$0,0.01,0.1,1,10$，以八折RMSE的算术均值选择，不由测试集选模。
岭目标为
\begin{equation}
\min_{\beta_0,\boldsymbol\beta}
\sum_i(y_i-\beta_0-\boldsymbol z_i^{\mathsf T}\boldsymbol\beta)^2
+\lambda\|\boldsymbol\beta\|_2^2,
\end{equation}
截距不惩罚。最终参数用完整训练集拟合，在首次测试TCAD提交前冻结模型、
标准化参数、Geant4测试输入、64条逐位预测及其哈希。
留出和外推完成后仅计算误差，不再按测试误差修改模型。

\subsection{分层验证与工程验收}
每个有效响应需满足完整求解到$4\ns$、有限数据、初态误差不超过$1\mV$、
预充误差不超过$10\,\mu\mathrm{V}$，并检查逐相KCL、四端口闭合和
后向欧拉右端点电荷积分。粒子源实际积分与目标Q的偏差门槛为5\%，
参考支路源项应为零。对物理双支路和离线参考同时比较两侧BL及差分，
防止两侧误差偶然抵消。

预登记代表点分别进行半时间步和0.7网格尺度检查，并为每个变体匹配自己的
对照和中性参考。$\Delta M$变化门槛为$0.5\mV$、$\Delta Q$为$0.01\fC$，
同时核对分类稳定性。训练交叉验证、数值细化和独立测试解决不同问题：
前者检验训练工况内的拟合选择，中者检查离散敏感性，后者检验未见尺寸条件。
任何一项通过都不能自动替代另外两项。
